\section{How format handlers reconstruct objects}
Native and legacy XML handling is organized into small element handlers, attribute setters and component-specific adapters. \class{SimpleSAX} delegates parser events to the active \class{ElementHandler}; this abstraction assumes an element contains text or child elements rather than arbitrary mixed content. It pools a bounded number of readers and clears handlers before returning a reader. The structure lets nested motor, component, configuration and result elements be interpreted by the classes that understand their semantics.

The native savers mirror the component vocabulary with type-specific property output. Matching import handlers interpret stored units and enumerations, select defaults for absent fields, reconstruct configuration IDs and references, and accumulate warnings. XML tag names and persisted identifiers are part of the format contract. A class rename or a new enum value can therefore require explicit compatibility handling even when the Java compiler accepts the change.

Storage also includes document preferences, custom materials and appearances. Attachments are resolved through an attachment factory and can originate from a local file or a ZIP entry. Keeping image bytes in an ORK archive makes a design's decal portable; a filename string alone would depend on another machine's directory layout. External FC definitions have the separately documented library/provenance contract and must not be assumed to behave like embedded decal files.

\sourceline{\source{core/src/main/java/info/openrocket/core/file/simplesax/SimpleSAX.java}{SimpleSAX}; \source{core/src/main/java/info/openrocket/core/file/simplesax/ElementHandler.java}{ElementHandler}; \source{core/src/main/java/info/openrocket/core/file/openrocket/importt/DocumentPreferencesHandler.java}{document preferences import}.}

\section{RockSim and RASAero interoperability}
RockSim import and export use an explicit mapping between RockSim entities and OpenRocket components. The import handlers interpret dimensions, materials, placement and other supported attributes. The export path builds JAXB data-transfer objects from a copied rocket, including the format's named stage structure, unit conversions and component-specific fields. That copy permits conversion adjustments without rewriting the open design. The presence of both directions does not guarantee a lossless round trip for all OpenRocket features.

RASAero import/export has a separate mapping for design geometry, launch conditions, recovery and supported simulation data. Its exporter also uses JAXB DTOs and collects warnings/errors when the source design does not fit the target representation. A target format's component types, allowed stage structures and aerodynamic assumptions constrain what can be expressed. The conversion layer, rather than the core component superclass, owns these compatibility rules.

Use the native ORK format to retain OpenRocket-specific design and simulation semantics. RockSim and RASAero files are interoperability products whose warnings and converted values need to travel with the result. In particular, a foreign application's simulation settings are not automatically an executable equivalent of MIT FC tasks, ensembles or active control surfaces.

\sourceline{\source{core/src/main/java/info/openrocket/core/file/rocksim/export/RockSimSaver.java}{RockSimSaver}; \source{core/src/main/java/info/openrocket/core/file/rasaero/export/RASAeroSaver.java}{RASAeroSaver}.}

\section{OBJ mesh export and geometry-only products}
The Wavefront OBJ subsystem creates mesh geometry from supported components and active instance transforms. \class{OBJExporterFactory} dispatches to body-tube, transition/nose, fin, tube-fin, lug, ring, mass-object, rail-button and motor exporters. Assemblies provide placement and selection structure; they are skipped as independent solid meshes while their selected children can be exported. Low-level cylinder, disk, tube and polygon builders create vertices, faces, normals and texture coordinates.

Options control selected components and children, repeated instances, motors, level of detail, coordinate transformation, scale, origin offset, triangulation, appearance, and separate component files. Material export writes an accompanying MTL representation. Mesh utilities manage groups, material assignments, vertex bounds and triangulation. Changing tessellation resolution changes the exported polygonal approximation, not the parametric source design or aerodynamic strip discretization.

OBJ is useful for external rendering and geometry workflows. It contains a snapshot of surfaces, not the OpenRocket tree's editable relationships, flight configurations, motor event rules or simulator. The fin editor also exports SVG outlines through \class{SVGBuilder}. It constructs a continuous fin-and-tab path and emits a separate tab path when a tab extends beyond the fin, with chosen stroke color and width. This is a fabrication/vector-drawing export, not an arbitrary-SVG-to-rocket modeling pipeline.

\sourceline{\source{core/src/main/java/info/openrocket/core/file/wavefrontobj/export/OBJExporterFactory.java}{OBJExporterFactory}; \source{core/src/main/java/info/openrocket/core/file/wavefrontobj/export/OBJExportOptions.java}{OBJExportOptions}; \source{core/src/main/java/info/openrocket/core/file/wavefrontobj/DefaultObj.java}{mesh representation}; \source{core/src/main/java/info/openrocket/core/file/svg/export/SVGBuilder.java}{SVGBuilder}.}

\section{Motor and component catalog ingestion}
RASP \code{.eng} and RockSim \code{.rse} motor readers turn curve samples and metadata into \class{ThrustCurveMotor} records. General/ZIP loaders dispatch among supported input forms. The motor database groups and indexes records for lookup and selection; curve identity, manufacturer aliases and digest information help distinguish records that have similar labels. The \code{thrustcurve} package handles online request/response parsing and downloads, while the GUI presents search and selection. These are data acquisition and resolution paths, not an alternative thrust integrator.

Component presets use typed keys and property maps, XML transfer objects, material definitions and catalog loaders. The many CSV loaders in \code{preset.loader} normalize differing catalog columns and units into the same preset model. Presets cover supported component families such as tubes, transitions, noses, fins, rings, launch hardware and recovery parts. Material databases distinguish bulk, surface and line densities before a component calculates mass.

Bundled data can be serialized or preprocessed by developer utilities to reduce startup work. Runtime loaders also support user-supplied data. A successful catalog import only makes a part or curve available; a document still needs to select or instantiate it. Once a preset is applied, editing the component can clear its catalog association while preserving the resulting dimensions.

\sourceline{\source{core/src/main/java/info/openrocket/core/database/motor/ThrustCurveMotorSet.java}{grouped motor records}; \source{core/src/main/java/info/openrocket/core/preset/ComponentPreset.java}{typed preset}; \source{core/src/main/java/info/openrocket/core/database/ComponentPresetDatabaseLoader.java}{preset loading}.}
